\documentclass[11pt]{article}

\usepackage[margin=1in]{geometry}
\usepackage{amsmath,amssymb,amsthm,mathtools,bm}
\usepackage{microtype}

\newtheorem{proposition}{Proposition}
\newtheorem{remark}{Remark}

\newcommand{\R}{\mathbb{R}}
\newcommand{\E}{\mathbb{E}}
\newcommand{\cP}{\mathcal{P}}
\newcommand{\cJ}{\mathcal{J}}
\newcommand{\cF}{\mathcal{F}}
\newcommand{\PiCouplings}{\Pi}
\newcommand{\norm}[1]{\left\lVert #1\right\rVert}
\newcommand{\tr}{\operatorname{tr}}
\newcommand{\Cov}{\operatorname{Cov}}
\newcommand{\diag}{\operatorname{diag}}
\newcommand{\argmin}{\operatorname*{arg\,min}}
\newcommand{\argmax}{\operatorname*{arg\,max}}
\newcommand{\equivmarg}{\equiv_{\mu,\nu}}

\title{Self-consistent optimal-transport formulations of two squared GW problems}
\author{}
\date{}

\begin{document}
\maketitle

\section*{High-level account}

Let
\[
\cJ(\pi)
=
\iint
\bigl(s_X(x,x')-s_Y(y,y')\bigr)^2
\,d\pi(x,y)\,d\pi(x',y')
\]
be the squared Gromov--Wasserstein objective over
$\pi\in\PiCouplings(\mu,\nu)$.
At a minimizer $\pi_\star$, the first variation of this quadratic
functional must be nonnegative in every feasible direction. Consequently,
$\pi_\star$ solves an ordinary Kantorovich problem whose cost is the
linearization of $\cJ$ at $\pi_\star$. Thus GW can be written as the
fixed-point condition
\[
\boxed{
\pi_\star\in
\operatorname{OT}_{c_{\Theta(\pi_\star)}}(\mu,\nu),
}
\]
where the cost parameter $\Theta(\pi_\star)$ is itself a moment matrix of
$\pi_\star$.

For the inner-product structures
\[
s_X(x,x')=x^\top x',
\qquad
s_Y(y,y')=y^\top y',
\]
the relevant matrix is
\[
M_\pi=\E_\pi[XY^\top],
\]
and an optimal coupling satisfies
\[
\pi_\star\in
\argmax_{\gamma\in\PiCouplings(\mu,\nu)}
\E_\gamma[X^\top M_{\pi_\star}Y].
\]
The full GW objective depends on $\pi$ only through $M_\pi$. Hence, when
both marginals are Gaussian, any optimal cross-moment can be realized by a
jointly Gaussian coupling. In equal dimension, with positive-definite
centered Gaussian marginals, every optimizer is in fact a deterministic
linear Gaussian coupling.

For squared Euclidean structures
\[
s_X(x,x')=\norm{x-x'}^2,
\qquad
s_Y(y,y')=\norm{y-y'}^2,
\]
one first lifts $x$ to $(\norm{x}^2,x,1)$. The same matrix fixed-point
statement then holds in the lifted space. For centered marginals it reduces
to
\[
\pi_\star\in
\argmax_{\gamma\in\PiCouplings(\mu,\nu)}
\E_\gamma\!\left[
\norm{X}^2\norm{Y}^2+4X^\top K_{\pi_\star}Y
\right],
\qquad
K_\pi=\E_\pi[XY^\top].
\]
Unlike the inner-product case, the radial fourth-order term
$\E[\norm{X}^2\norm{Y}^2]$ is not determined by the cross-covariance.
This permits non-Gaussian dependence to outperform every Gaussian
coupling. An explicit example is given below for
$\mathcal N(0,1)$ and $\mathcal N(0,I_{10})$.

\section{The general linearization principle}

Let $\mu\in\cP(\R^d)$ and $\nu\in\cP(\R^e)$ have enough moments for all
expressions below to be finite. Let $s_X$ and $s_Y$ be symmetric structure
functions, and define
\[
L\bigl((x,y),(x',y')\bigr)
=
\bigl(s_X(x,x')-s_Y(y,y')\bigr)^2.
\]
For $\pi\in\PiCouplings(\mu,\nu)$, set
\[
\cJ(\pi)=\iint L(z,z')\,d\pi(z)\,d\pi(z'),
\qquad
\ell_\pi(z)=\int L(z,z')\,d\pi(z').
\]

\begin{proposition}[Self-consistent OT condition]
If $\pi_\star$ minimizes $\cJ$ over $\PiCouplings(\mu,\nu)$, then
\[
\pi_\star
\in
\argmin_{\gamma\in\PiCouplings(\mu,\nu)}
\int \ell_{\pi_\star}(x,y)\,d\gamma(x,y).
\]
Thus every GW minimizer is an OT minimizer for the cost
$c_{\pi_\star}=\ell_{\pi_\star}$ generated by that same coupling.
\end{proposition}

\begin{proof}
For any $\gamma\in\PiCouplings(\mu,\nu)$, put
\[
\pi_t=(1-t)\pi_\star+t\gamma.
\]
Since $\pi_\star$ is a minimizer,
\[
0\leq
\left.\frac{d}{dt}\cJ(\pi_t)\right|_{t=0^+}
=
2\int \ell_{\pi_\star}(z)\,d(\gamma-\pi_\star)(z).
\]
Rearranging gives the claimed Kantorovich optimality condition.
\end{proof}

The condition is necessary; by itself it is not generally sufficient for
global GW optimality, because it is a nonlinear fixed-point condition.

\section{Bilinear feature kernels and the matrix parameter}

Suppose that the two structures admit finite-dimensional bilinear
representations
\[
s_X(x,x')=\phi(x)^\top A\phi(x'),
\qquad
s_Y(y,y')=\psi(y)^\top B\psi(y'),
\]
where $A$ and $B$ are symmetric matrices. Define the cross-feature moment
\[
M_\pi=\int \phi(x)\psi(y)^\top\,d\pi(x,y).
\]
Independence of the two copies drawn from $\pi$ gives
\[
\iint s_X(x,x')s_Y(y,y')\,d\pi\,d\pi
=
\tr\!\left(A M_\pi B M_\pi^\top\right).
\]
Therefore
\[
\cJ(\pi)
=
C_{\mu,\nu}
-2\tr\!\left(A M_\pi B M_\pi^\top\right),
\]
where $C_{\mu,\nu}$ depends only on the marginals.

Moreover, after deleting a function of $x$ and a function of $y$, which do
not affect a Kantorovich minimizer,
\[
\ell_\pi(x,y)
\equivmarg
-2\phi(x)^\top A M_\pi B\psi(y).
\]
Here $c\equivmarg\widetilde c$ means
\[
c(x,y)-\widetilde c(x,y)=f(x)+g(y)
\]
for some $f$ and $g$. Such terms have constant integral over
$\PiCouplings(\mu,\nu)$.

Consequently, with
\[
\Theta_\pi=A M_\pi B,
\]
every GW optimizer obeys
\[
\boxed{
\pi_\star\in
\argmin_{\gamma\in\PiCouplings(\mu,\nu)}
\int -\phi(x)^\top\Theta_{\pi_\star}\psi(y)\,d\gamma(x,y).
}
\]
This is the desired ordinary OT problem with a matrix-valued parameter
that is determined self-consistently by the optimal coupling.

\section{Inner-product GW and Gaussian couplings}

Take
\[
s_X(x,x')=x^\top x',
\qquad
s_Y(y,y')=y^\top y'.
\]
Then $\phi(x)=x$, $\psi(y)=y$, $A=I_d$, $B=I_e$, and
\[
M_\pi=\E_\pi[XY^\top].
\]
The objective and its self-consistent OT formulation are
\begin{align*}
\cJ_{\mathrm{ip}}(\pi)
&=C_{\mu,\nu}-2\norm{M_\pi}_{\mathrm F}^2,\\
\pi_\star
&\in
\argmax_{\gamma\in\PiCouplings(\mu,\nu)}
\int x^\top M_{\pi_\star}y\,d\gamma(x,y).
\end{align*}
Thus inner-product GW maximizes the squared Frobenius norm of a second
cross-moment.

\subsection{Why a Gaussian optimizer exists}

Let
\[
\mu=\mathcal N(m_0,\Sigma_0),
\qquad
\nu=\mathcal N(m_1,\Sigma_1).
\]
For any coupling $\pi$, let
\[
K_\pi=\Cov_\pi(X,Y),
\qquad
M_\pi=m_0m_1^\top+K_\pi.
\]
Because $K_\pi$ comes from an actual coupling, the block covariance matrix
\[
\Sigma_\pi=
\begin{pmatrix}
\Sigma_0 & K_\pi\\
K_\pi^\top & \Sigma_1
\end{pmatrix}
\]
is positive semidefinite. Hence the jointly Gaussian law
\[
\widehat\pi
=
\mathcal N\!\left(
\binom{m_0}{m_1},
\Sigma_\pi
\right)
\]
is a coupling of $\mu$ and $\nu$ with exactly the same $M_\pi$.

Since $\cJ_{\mathrm{ip}}$ depends on $\pi$ only through $M_\pi$,
\[
\cJ_{\mathrm{ip}}(\widehat\pi)
=
\cJ_{\mathrm{ip}}(\pi).
\]
Applying this replacement to any optimizer proves that a jointly Gaussian
inner-product GW optimizer always exists.

The fixed-point OT formula gives the same conclusion more geometrically.
If $M_{\pi_\star}$ is invertible, set
\[
z=M_{\pi_\star}^\top x.
\]
Then maximizing
\[
x^\top M_{\pi_\star}y=z^\top y
\]
is equivalent, up to marginal terms, to minimizing $\norm{z-y}^2$, since
\[
-z^\top y
=
\frac12\norm{z-y}^2
-\frac12\norm{z}^2
-\frac12\norm{y}^2.
\]
The quadratic OT map between two nondegenerate Gaussian laws is affine.
Composing it with $x\mapsto M_{\pi_\star}^\top x$ makes $Y$ an affine
function of $X$, so the optimal joint law is Gaussian.

\subsection{Equal-dimensional centered case: every optimizer is Gaussian}

Assume now
\[
\mu=\mathcal N(0,\Sigma_0),
\qquad
\nu=\mathcal N(0,\Sigma_1),
\qquad
\Sigma_0,\Sigma_1\succ0,
\]
with both measures on $\R^d$. A matrix $K=\E[XY^\top]$ is feasible exactly
when it can be written as
\[
K=\Sigma_0^{1/2}C\Sigma_1^{1/2},
\qquad
\norm{C}_{\mathrm{op}}\leq1.
\]
Hence the moment problem is
\[
\max_{\norm{C}_{\mathrm{op}}\leq1}
\tr\!\left(\Sigma_0 C\Sigma_1 C^\top\right).
\]

Let the eigenvalues of $\Sigma_0$ and $\Sigma_1$ be
\[
\alpha_1\geq\cdots\geq\alpha_d>0,
\qquad
\beta_1\geq\cdots\geq\beta_d>0.
\]
The trace rearrangement inequality gives
\[
\max_{\norm{C}_{\mathrm{op}}\leq1}
\tr\!\left(\Sigma_0 C\Sigma_1 C^\top\right)
=
\sum_{i=1}^d\alpha_i\beta_i.
\]
The maximum is attained by an orthogonal $C_\star$ aligning the ordered
eigenspaces, with arbitrary signs inside one-dimensional eigenspaces and
arbitrary orthogonal choices inside repeated eigenspaces. Since both
covariances are positive definite, every maximizing contraction saturates
all directions and is orthogonal on $\R^d$.

For such a maximizer,
\[
K_\star=\Sigma_0^{1/2}C_\star\Sigma_1^{1/2}
\]
and
\[
\Sigma_1-K_\star^\top\Sigma_0^{-1}K_\star
=
\Sigma_1^{1/2}(I-C_\star^\top C_\star)\Sigma_1^{1/2}
=0.
\]
Therefore every coupling having this optimal cross-covariance satisfies
\[
Y=K_\star^\top\Sigma_0^{-1}X
\qquad\text{almost surely}.
\]
It is a deterministic linear image of a Gaussian variable and is therefore
a jointly Gaussian coupling. Thus, in this nondegenerate equal-dimensional
centered setting, not only does a Gaussian optimizer exist: every optimizer
is Gaussian.

\section{Squared Euclidean structures via a quadratic lift}

Now take
\[
s_X(x,x')=\norm{x-x'}^2,
\qquad
s_Y(y,y')=\norm{y-y'}^2.
\]
Introduce
\[
\Phi_d(x)=
\begin{pmatrix}
\norm{x}^2\\
x\\
1
\end{pmatrix},
\qquad
H_d=
\begin{pmatrix}
0 & 0 & 1\\
0 & -2I_d & 0\\
1 & 0 & 0
\end{pmatrix}.
\]
Then
\[
\norm{x-x'}^2
=
\Phi_d(x)^\top H_d\Phi_d(x').
\]
Likewise define $\Phi_e$ and $H_e$ for $y$. With the lifted cross-moment
\[
\widetilde M_\pi
=
\E_\pi\!\left[\Phi_d(X)\Phi_e(Y)^\top\right],
\]
the general bilinear-feature calculation yields
\[
\Theta_\pi^{\mathrm{dist}}
=
H_d\widetilde M_\pi H_e
\]
and
\[
\boxed{
\pi_\star\in
\argmin_{\gamma\in\PiCouplings(\mu,\nu)}
\int
-\Phi_d(x)^\top
\Theta_{\pi_\star}^{\mathrm{dist}}
\Phi_e(y)
\,d\gamma(x,y).
}
\]
Thus squared-distance GW also has a matrix-parameterized OT fixed point,
but the matrix acts on quadratic lifts rather than on $x$ and $y$ alone.

\subsection{Centered simplification}

Assume $\E[X]=\E[Y]=0$ and write
\[
K_\pi=\E_\pi[XY^\top].
\]
For fixed $(x,y)$, define
\[
h_\pi(x,y)
=
\E_\pi\!\left[
\norm{x-X'}^2\norm{y-Y'}^2
\right].
\]
Expanding the product and deleting marginal-only terms gives
\[
h_\pi(x,y)
\equivmarg
\norm{x}^2\norm{y}^2+4x^\top K_\pi y.
\]
Since the linearized squared-discrepancy cost is
\[
f(x)+g(y)-2h_\pi(x,y),
\]
every optimizer satisfies
\[
\boxed{
\pi_\star\in
\argmax_{\gamma\in\PiCouplings(\mu,\nu)}
\int
\left(
\norm{x}^2\norm{y}^2+4x^\top K_{\pi_\star}y
\right)
\,d\gamma(x,y).
}
\]
This is an ordinary OT problem with polynomial cost
\[
c_{K_\star}(x,y)
=
-\norm{x}^2\norm{y}^2-4x^\top K_\star y,
\qquad
K_\star=\E_{\pi_\star}[XY^\top].
\]

There is also a useful global form. Let $(X',Y')$ be an independent copy of
$(X,Y)$ and set
\[
a=\E\norm{X}^2,
\qquad
b=\E\norm{Y}^2,
\qquad
q_\pi=\E_\pi[\norm{X}^2\norm{Y}^2].
\]
A direct expansion gives
\[
\E\!\left[
\norm{X-X'}^2\norm{Y-Y'}^2
\right]
=
2q_\pi+2ab+4\norm{K_\pi}_{\mathrm F}^2.
\]
The other terms in the squared GW discrepancy depend only on the
marginals. Therefore minimizing squared-distance GW is equivalent to
maximizing
\[
\boxed{
\cF(\pi)
=
q_\pi+2\norm{K_\pi}_{\mathrm F}^2.
}
\]
The new quantity $q_\pi$ is a mixed fourth-order radial moment. This is the
term that breaks Gaussian closure.

\section{Counterexample: Gaussian marginals, but no Gaussian optimizer}

Let
\[
\mu=\mathcal N(0,1)
\quad\text{on }\R,
\qquad
\nu=\mathcal N(0,I_{10})
\quad\text{on }\R^{10}.
\]

\subsection{Best value attainable by a jointly Gaussian coupling}

For a jointly Gaussian coupling, write
\[
k=\E[XY^\top]\in\R^{1\times10}.
\]
Positive semidefiniteness of the block covariance gives
\[
\norm{k}^2\leq1.
\]
By Isserlis' formula,
\[
q_\pi
=
\E[X^2\norm{Y}^2]
=
\sum_{j=1}^{10}\E[X^2Y_j^2]
=
10+2\norm{k}^2.
\]
Consequently every jointly Gaussian coupling satisfies
\[
\cF(\pi)
=
q_\pi+2\norm{k}^2
=
10+4\norm{k}^2
\leq14.
\]
The value $14$ is attained, for example, by taking $Y_1=X$ and making
$Y_2,\ldots,Y_{10}$ independent standard Gaussians.

\subsection{A better non-Gaussian coupling}

Draw
\[
Y\sim\mathcal N(0,I_{10})
\]
and put
\[
S=\norm{Y}^2\sim\chi^2_{10}.
\]
Let $F_{10}$ and $F_1$ denote the distribution functions of
$\chi^2_{10}$ and $\chi^2_1$. Define
\[
U=F_{10}(S),
\qquad
A=F_1^{-1}(U).
\]
Then $U$ is uniform on $(0,1)$ and $A\sim\chi^2_1$. Let $\varepsilon$ be an
independent Rademacher variable and set
\[
X=\varepsilon\sqrt{A}.
\]
It follows that $X\sim\mathcal N(0,1)$, so this is a valid coupling of the
two Gaussian marginals. Moreover,
\[
\E[X\mid Y]=0,
\qquad
k=\E[XY^\top]=0.
\]
On the other hand, the squared radii are comonotonically coupled, and hence
\[
q_\pi
=
\E[AS]
=
\int_0^1
F_1^{-1}(u)F_{10}^{-1}(u)\,du
=
15.85744810\ldots.
\]
Thus
\[
\cF(\pi)=15.85744810\ldots>14.
\]
Because larger $\cF$ means smaller squared-distance GW objective, this
non-Gaussian coupling is strictly better than every jointly Gaussian
coupling. Therefore no Gaussian coupling can be optimal for this pair of
Gaussian marginals.

The joint law is non-Gaussian: $X^2$ is a deterministic increasing function
of $\norm{Y}^2$, while
\[
\Cov(X,Y)=0.
\]
Under joint Gaussianity, zero cross-covariance would imply independence,
which is impossible here.

\subsection{Why the phenomenon is not a numerical accident}

For general $n$, define
\[
I_{1,n}
=
\int_0^1F_1^{-1}(u)F_n^{-1}(u)\,du.
\]
The standardized chi-square quantiles satisfy, in $L^2(0,1)$,
\[
\frac{F_n^{-1}(u)-n}{\sqrt{2n}}
\longrightarrow
\Phi^{-1}(u),
\]
where $\Phi$ is the standard normal distribution function. Hence
\[
I_{1,n}
=
n+\sqrt{2n}\,c+o(\sqrt n),
\qquad
c=
\int_0^1F_1^{-1}(u)\Phi^{-1}(u)\,du>0.
\]
The positivity follows because both quantile functions are strictly
increasing and nonconstant. Therefore
\[
I_{1,n}>n+4
\]
for all sufficiently large $n$, while the best jointly Gaussian value is
exactly $n+4$. The counterexample is therefore structural: the radial
mixed fourth moment can gain an order-$\sqrt n$ amount that is invisible
to Gaussian cross-covariance optimization.

\begin{remark}[Same ambient dimension and nondegenerate covariances]
The preceding example uses different intrinsic dimensions, which is
natural for GW. It can also be converted into a same-dimensional full-rank
example. On $\R^{10}$, take
\[
\mu=\mathcal N(0,I_{10}),
\qquad
\nu_\delta
=
\mathcal N\!\left(
0,\diag(1,\delta,\ldots,\delta)
\right).
\]
For every jointly Gaussian coupling,
\[
\cF(\pi)\leq14(1+9\delta).
\]
Indeed,
\[
\tr(\Cov(Y))=1+9\delta
\]
and the maximal squared Frobenius norm of a feasible cross-covariance is
also $1+9\delta$.

For a non-Gaussian competitor, couple
$\norm{X}^2\sim\chi^2_{10}$ comonotonically with $\norm{Y}^2$, and choose
the direction of $X$ independently and uniformly on the sphere. This gives
$K_\pi=0$. As $\delta\downarrow0$, its value of $\cF$ converges to
\[
I_{1,10}=15.85744810\ldots,
\]
whereas the Gaussian upper bound converges to $14$. Hence the strict
inequality persists for all sufficiently small $\delta>0$. Both Gaussian
marginals can therefore be chosen nondegenerate and supported on the same
Euclidean space.
\end{remark}

\section*{Conclusion}

Both squared GW problems admit the same formal description:
\[
\pi_\star
\in
\operatorname{OT}_{c_{\Theta(\pi_\star)}}(\mu,\nu).
\]
For inner products,
\[
\Theta(\pi)=\E[XY^\top],
\]
and the GW objective sees only a second cross-moment. Gaussian marginals
are therefore closed under optimal Gaussian replacement, and in the
nondegenerate equal-dimensional centered case every optimizer is linear
Gaussian.

For squared Euclidean structures, the corresponding matrix acts on the
lift
\[
(\norm{x}^2,x,1).
\]
Equivalently, in the centered case the fixed-point OT cost contains both a
bilinear term and the radial product
\[
-\norm{x}^2\norm{y}^2.
\]
That fourth-order term can encode nonlinear radial dependence while
keeping both marginals Gaussian, which is exactly why the Gaussian
conclusion fails.

\end{document}